\section{Mixed boundary characters and the collision extension}
\label{sec:mixed-boundary-collision-extension}

This calculation applies the existing collision differential extension
to the computed characters of the all-prime mixed boundary module.
The two operators retain their distinct origins: the boundary operator
is an idele dilation; the collision operator is a specified Frobenius
lift. We construct their exact comparison complex and its extension
classes. No identification of these origins, arithmetic purity, or
sign of the original Weil pairing is presumed.

\subsection{The original boundary and its integral coefficient maps}

Let \(\mathcal P\) be all rational primes, \(r\ge2\), \(1\le m<r\),
and \(j\ge1\). Retain
\begin{equation}\tag{BCE1}
\begin{gathered}
\mathcal A_{r,m}=\{A\subset\{1,\ldots,r\}:|A|=m\},\qquad
X_{r,m}=\mathcal A_{r,m}^{\{\infty\}\sqcup\mathcal P},\\
E_{r,m}=\{(A,A,\ldots):A\in\mathcal A_{r,m}\},\\
I_{r,m}(R)=\{f\in\operatorname{LC}(X_{r,m},R):
                              f|_{E_{r,m}}=0\},\\
B_{r,m,j}(R)=I_{r,m}(R)\otimes_R
                   \Lambda^j\left(\bigoplus_pR\vartheta_p\right).
\end{gathered}
\end{equation}
Locally constant means depending on finitely many place coordinates;
the discrete coefficient ring is indicated. These coordinates are
evaluation/integration choices at each place, not a replacement for
the coefficient-support lattice \(L\).

The original GSB8, GSB15 and GSB16 calculation gives an injection
\begin{equation}\tag{BCE2}
\partial^\theta_{r,j}:B_{r,m,j}(\mathbb C)
 \longrightarrow H_{j-1}(\mathbb Q^\times,K_r^\theta),
\qquad
U_a\partial^\theta_{r,j}(b)
=\left(|a_\infty|_\infty\prod_p|a_p|_p\right)^m
                                  \partial^\theta_{r,j}(b).
\end{equation}
Here the original Schwartz source is \(T=\mathcal S(\mathbb A_\mathbb Q)\),
\(E\xi(x)=\sum_{q\in\mathbb Q^\times}\xi(qx)\), and \(b:T\to
\operatorname{LC}(\{0,1\}^{\{\infty\}\sqcup\mathcal P},\mathbb C)\)
tensors evaluation at zero and additive integration. With algebraic
tensors, \(K_r^\theta=\ker b^{\otimes r}\cap\ker E^{\otimes r}\).
The source of this connecting map is the group homology of the kernel
of the endpoint evaluations. The complete proof, including topology,
all-prime representatives and injectivity, is
\href{https://github.com/KokunoYumeto/zeta-function-research-reader/blob/7bb431df02b42681b3c7335fbf272a3771f777ef/workbenches/splitzero-tandem/research-self-audit/supporting_proofs/GLOBAL_SECONDARY_BOUNDARY_EVALUATION.tex#L1}
{GSB1--22a}. Its original adelic conventions are those of Connes and
Connes--Consani--Marcolli \cite{BCEConnes,BCECCM}. BCE2 is an input from that
proof, not a new assertion of integral Schwartz exactness.

Put \(R_0=\mathbb Z_{(3)}\). For every coefficient ring \(R\) over
\(R_0\), there is an exact scalar-extension identification
\begin{equation}\tag{BCE3}
I_{r,m}(R_0)\otimes_{R_0}R=I_{r,m}(R),\qquad
B_{r,m,j}(R_0)\otimes_{R_0}R=B_{r,m,j}(R).
\end{equation}
Here both source modules over \(R_0\) are free. To prove all these
claims, order the primes and use successively the first \(n\).
At each finite stage, a basis for the kernel of endpoint restriction
consists of the indicator functions of patterns other than the
constant patterns. Restriction to constant patterns is exactly
endpoint evaluation, so no other relation occurs. Refinement sends
each old basis vector to the sum of its disjoint child patterns.
Choosing one child of each old nonconstant pattern gives an explicit
linear retraction. Each such block quotient is free by eliminating
that one child, and new nonconstant children of old constant patterns
are additional independent basis vectors. Thus each stage splits as
the preceding stage plus a finite free complement. Their union is
free. This proof works over every \(R\) and proves the first equality
in BCE3. The exterior factor has the basis of increasing distinct
prime wedges; tensoring these bases proves the second equality.

Consequently the following are actual coefficient arrows:
\begin{equation}\tag{BCE4}
B_{r,m,j}(R_0)\longrightarrow B_{r,m,j}(\mathbb C)
\xrightarrow{\partial^\theta_{r,j}}\operatorname{im}\partial^\theta_{r,j},
\qquad
B_{r,m,j}(R_0)\longrightarrow B_{r,m,j}(\mathbb Z_3).
\end{equation}
Both first coefficient maps are injective by the free bases just
proved. There is no arrow \(\mathbb C\to\mathbb Q_3\) in BCE4.

Choose the actual idele
\(\alpha=(3,(1)_p)\). Its modulus is \(3\), so its action in BCE2
is \(a_m=3^m\). This is \emph{not} the diagonal rational idele \(3\),
whose modulus is \(3\cdot3^{-1}=1\). The lattice in BCE4 is stable
under \(U_\alpha\) and its nonnegative powers. Its action is not
surjective there; the inverse requires inverting 3. The entire
complex idele representation in BCE2 is retained, and is not
replaced by this integral semigroup lattice.

\subsection{The retained collision coefficients}

Write \(O=\mathbb Z_3\), \(K=\mathbb Q_3\), retain \(c\in O\), and put
\begin{equation}\tag{BCE5}
\begin{gathered}
C=O[\epsilon,T]/(\epsilon^2,T^2-6\epsilon),\qquad
z=\epsilon T,\quad M=O\epsilon\oplus OT\oplus Oz,\\
F_M(T)=9cz,\quad F_M(\epsilon)=F_M(z)=0 .
\end{gathered}
\end{equation}
The original algebra has basis \(1,\epsilon,T,z\), by successive
monic division. In particular \(F_M^2=0\); the algebra map is the
identity on \(O\). These are the original collision coefficients of
CFF1--4, not a Frobenius assigned to BCE2.

The cotangent differential has presentation
\[
C^2\xrightarrow{\left(\begin{smallmatrix}2\epsilon&-6\\0&2T
\end{smallmatrix}\right)} C\,d\epsilon\oplus C\,dT .
\]
It yields the exact \(O\)-module and extension
\begin{equation}\tag{BCE6}
\begin{gathered}
\Omega=Ou\oplus Ov\oplus Oq\oplus(O/3)w,\\
u=d\epsilon,\quad v=dT,\quad q=\epsilon dT,\quad
w=T\,d\epsilon-2\epsilon dT,\\
0\longrightarrow M\xrightarrow{d_M}\Omega
\xrightarrow{\rho}Q=O/9\longrightarrow0,\\
d_M(\epsilon)=u,\quad d_M(T)=v,\quad d_M(z)=w+3q,\\
\rho(u)=\rho(v)=0,\quad\rho(q)=1,\quad\rho(w)=-3,\qquad
9q=3d_M(z).
\end{gathered}
\end{equation}
For verification directly from the eight generators of
\(C\,d\epsilon\oplus C\,dT\), the relation columns and their
multiples impose
\(\epsilon u=zu=0,\ Tv=3u,\ zv=0,\ 3(Tu-2\epsilon v)=0\).
All remaining multiples follow from these. Eliminating the four
zero or dependent generators and setting \(Tu=w+2q\) gives exactly
the displayed three free generators and the order-three generator.
The image of \(a\epsilon+bT+tz\) has free coordinates \(a,b,3t\);
it is zero only for \(a=b=t=0\). In its cokernel, \(u=v=0\),
\(w=-3q\) and \(3w=0\), giving precisely \(O/9\).
Changing the lift \(q\) changes \(3z\) by \(9M\).
Thus its extension class is \([3z]\in M/9M\), of exact order three.

The actual induced operators on this sequence are
\begin{equation}\tag{BCE7}
F_\Omega(u)=F_\Omega(q)=F_\Omega(w)=0,\qquad
F_\Omega(v)=27cq,\qquad F_Q=0,\qquad F_\Omega d_M=d_MF_M .
\end{equation}
Indeed \(d(9cz)=9c(w+3q)=27cq\), since \(3w=0\).
All operators square to zero. The coefficients and formulas in
BCE5--7 are already defined over \(\mathbb Z_{(3)}[c]\).
They can be specialized separately to a chosen \(c\in O\) and to a
chosen complex parameter. Such specializations are not a map from
one coefficient field to the other.

\subsection{The exact operator comparison and all tensor orders}

For an \(O\)-module \(E\) with endomorphism \(F_E\), define
\begin{equation}\tag{BCE8}
\mathcal K_m(E)=
[\,E\xrightarrow{\,L_{m,E}=3^mI-F_E\,}E\,]
\quad\hbox{in cohomological degrees }0,1 .
\end{equation}
This is a specified comparison object between the two endomorphisms.
Its degree-zero cohomology is the actual intertwiner module
\(\operatorname{Hom}_{O[t]}(O_{3^m},E_F)\), where \(t\) acts on the
source as \(3^m\) and on the target as \(F_E\). Its degree-one
cohomology is \(\operatorname{Ext}^1_{O[t]}(O_{3^m},E_F)\), and
the higher Ext groups are zero. Here is the proof, including signs.
There is an exact free resolution
\begin{equation}\tag{BCE9}
0\longrightarrow O[t]\xrightarrow{\,3^m-t\,}O[t]
\longrightarrow O_{3^m}\longrightarrow0 .
\end{equation}
Division by the monic polynomial \(t-3^m\) proves the cokernel and
the leading term proves injectivity. Applying \(\operatorname{Hom}(-,E)\)
gives BCE8 with the displayed sign.
More concretely, a class \(x\in E/(3^m-F_E)E\) gives the extension
\[
E\oplus O\widetilde 1,\qquad
t\widetilde1=3^m\widetilde1-x,\qquad t|_E=F_E .
\]
Every extension has this form because its underlying \(O\)-module
extension splits. Replacing \(\widetilde1\) by
\(\widetilde1+y\) changes \(x\) to \(x+(3^m-F_E)y\).
This proves the classification and its splitting criterion without
omitting the integral lattice.

For \(k\ge1\) put \(M_k=M^{\otimes_O k}\), with
\(F_k=F_M^{\otimes k}\), on the original words in \(\epsilon,T,z\).
Let \(n=3^k\), \(a=3^m\), and \(b=(9c)^k\).
The map \(F_k\) sends \(T^{\otimes k}\) to \(b z^{\otimes k}\)
and kills every other word. Thus
\begin{equation}\tag{BCE10}
L_{m,k}=aI-F_k,\qquad
L_{m,k}^{-1}=a^{-1}I+a^{-2}F_k\quad\hbox{over }K .
\end{equation}
Multiplying the two displayed operators gives \(I\), since
\(F_k^2=0\). In particular \(\ker L_{m,k}=0\) on the original
free \(O\)-module, and BCE10 is a contracting homotopy for its
rationalized two-term complex, not a deletion of its integral
cokernel.

Retain \(v_3(0)=+\infty\), and write
\(\nu=\min\{m,2k+kv_3(c)\}\), interpreting the second term as
\(+\infty\) when \(c=0\). The full answer is
\begin{equation}\tag{BCE11}
H^1\mathcal K_m(M_k)\cong
(O/3^m)^{\,3^k-2}\oplus O/3^\nu\oplus O/3^{\,2m-\nu}.
\end{equation}
The two exceptional coordinates have the retained matrix
\(\left(\begin{smallmatrix}a&-b\\0&a\end{smallmatrix}\right)\)
in the order \(z^{\otimes k},T^{\otimes k}\).
If \(a\mid b\), add \(b/a\) times the first column to the second,
obtaining \(aI\). Otherwise \(b=3^\nu u\) with \(u\) a unit and
\(\nu<m\). Move the entry \(-b\) to the first diagonal position,
multiply that row by its unit inverse, and use the fact that it
divides all entries to eliminate its row and column. The remaining
diagonal entry has valuation \(2m-\nu\), since the original
determinant is \(a^2\). All these changes are invertible over \(O\).
The other \(3^k-2\) coordinates are multiplication by \(a\).
This proves BCE11, including \(m=0\), \(c=0\) and all tensor orders.
The row and column maps do not replace the original coordinates
when using the connecting maps below.

There is also a perfect torsion pairing, with the transpose
coefficient object retained:
\begin{equation}\tag{BCE12}
\operatorname{coker}L_{m,k}\times\operatorname{coker}L_{m,k}^{t}
\longrightarrow K/O,\qquad
([x],[y])\longmapsto y^tL_{m,k}^{-1}x\pmod O .
\end{equation}
Changing \(x\) by \(L_{m,k}u\) adds \(y^tu\in O\);
changing \(y\) by \(L_{m,k}^tv\) adds \(v^tx\in O\).
If a class pairs to zero against all integral \(y\), every
coordinate of \(L_{m,k}^{-1}x\) is integral, so \(x\) is in
\(L_{m,k}M_k\). The same argument with the transpose proves the
other radical is zero. BCE11 gives equal finite orders on both
sides, and their characters into \(K/O\) have those same orders:
for a summand \(O/3^h\), its character is determined by an element
of \(3^{-h}O/O\), also of order \(3^h\). Thus the pairing is perfect.
It is a torsion-valued pairing, not a real ordered or Hermitian form.

\subsection{Applying the order-nine extension: zero and nonzero images}

By BCE7, the original exact sequence BCE6 gives a short exact
sequence of the actual comparison complexes:
\begin{equation}\tag{BCE13}
0\longrightarrow\mathcal K_m(M)
\xrightarrow{d_M}\mathcal K_m(\Omega)
\xrightarrow{\rho}\mathcal K_m(Q)\longrightarrow0 .
\end{equation}
For \(m\ge1\), the complete kernels needed for the connecting map are
\begin{equation}\tag{BCE14}
H^0\mathcal K_m(M)=0,\quad
H^0\mathcal K_m(\Omega)=(O/3)w,\quad
H^0\mathcal K_m(Q)=Q[3^m]
      =3^{\max(0,2-m)}O/9 .
\end{equation}
To check the middle equality, write \(xu+yv+hq+tw\).
Its image has free coordinates \(3^mx,\ 3^my,\ 3^mh-27cy\);
all vanish only for \(x=y=h=0\). Multiplication by \(3^m\)
kills \(w\). For \(m=0\) all three comparison maps are invertible
and every cohomology group is zero.

Let \(s=\max(0,2-m)\), and take \(3^s\) as the displayed generator
of the last module of BCE14. Its lift in \(\Omega\) is \(3^sq\).
Its differential is
\[
(3^m-F_\Omega)(3^sq)=3^{m+s}q
                  =3^{m+s-1}d_M(z).
\]
The sign convention BCE8 therefore gives
\begin{equation}\tag{BCE15}
\delta_m([3^s])=[3^{m+s-1}z]
             \in M/(3^m-F_M)M .
\end{equation}
The class \([z]\) has exact order \(3^m\): if a multiple of \(z\)
is \(L_{m,M}(x\epsilon+yT+hz)\), its \(T\) and \(\epsilon\)
coefficients force \(x=y=0\), and its \(z\) coefficient is
\(3^mh\). Consequently
\begin{equation}\tag{BCE16}
\delta_1=0,\qquad
\delta_m(1)=[3^{m-1}z]\ne0,\quad
3\delta_m(1)=0,\quad
\ker\delta_m=3O/9 \quad(m\ge2).
\end{equation}
This is independent of the retained parameter \(c\). The kernel
statement also matches the preceding map in the exact sequence:
the actual map \((O/3)w\to Q[3^m]\) sends \(w\) to \(-3\).

For clarity the full first mixed case \(m=1\) is
\begin{equation}\tag{BCE17}
\begin{gathered}
H^1\mathcal K_1(M)=(O/3)\{\epsilon,T,z\},\qquad
H^1\mathcal K_1(\Omega)=(O/3)\{u,v,q,w\},\\
H^1\mathcal K_1(Q)=O/3,\qquad
d_{M,*}(\epsilon)=u,\quad d_{M,*}(T)=v,\quad d_{M,*}(z)=w,\\
\rho_*q=1,\quad \rho_*u=\rho_*v=\rho_*w=0 .
\end{gathered}
\end{equation}
The free image in \(\Omega\) is exactly \(3Ou+3Ov+3Oq\):
the off-diagonal \(27cq\) lies in it, and no \(w\) lies in the
image. This proves BCE17 and its zero connecting map directly,
rather than by forgetting the original order-nine quotient.
For \(m\ge1\) in general the same argument and two-by-two
elimination give
\begin{equation}\tag{BCE18}
\begin{gathered}
H^1\mathcal K_m(\Omega)\cong
O/3^m\oplus O/3^\mu\oplus O/3^{2m-\mu}\oplus O/3,\\
\mu=\min\{m,3+v_3(c)\},\qquad
H^1\mathcal K_m(Q)=O/3^{\min(m,2)} .
\end{gathered}
\end{equation}
The original map is still BCE6; BCE18 does not authorize replacing
it by an unspecified map between cyclic presentations.

The \(z\)-coordinate functional in BCE12 is especially explicit:
\begin{equation}\tag{BCE19}
\lambda_m([x_\epsilon\epsilon+x_TT+x_zz])
=\frac{x_z}{3^m}+\frac{9c\,x_T}{3^{2m}}\pmod O .
\end{equation}
Its value on \(L_{m,M}y\) is \(y_z\in O\), so it is well-defined.
It detects the whole displayed connecting generator:
\begin{equation}\tag{BCE20}
\lambda_m(\delta_m(1))=\frac13\pmod O\ne0
\qquad(m\ge2).
\end{equation}
Thus rational acyclicity in BCE10 is compatible with an integral
extension and its explicit nonzero torsion character. An additive
map from this order-three image to a characteristic-zero vector
space is zero, since its target has no element killed by 3.
This assertion concerns that specified additive map, not all
possible arithmetic uses of the extension.

\subsection{The existing cotangent injection and its exact repair}

Retain \(H=Ov_1\oplus Ov_2\oplus Ov_3=H^{-1}(L_{C/O})\).
The original cotangent cycles are
\[
v_1=(\epsilon,0),\quad v_2=(z,0),\quad v_3=(3T,z).
\]
They exhaust the kernel of the matrix preceding BCE6: its lower
equation forces its second coordinate to be a multiple of \(z\),
and its upper equation then fixes the \(T\)-coefficient of its
first coordinate. The actual Frobenius on \(H\) is zero.
Indeed the polynomial lift sends its two relation generators
\(\epsilon^2,T^2-6\epsilon\) to \(0,81c^2T^2\epsilon^2\).
The induced degree-minus-one semilinear matrix is
\(\left(\begin{smallmatrix}0&486c^2\epsilon\\0&0\end{smallmatrix}\right)\);
the second coordinates of the displayed cycles are \(0,0,z\),
whose Frobenius values all vanish.

The retained injection and defect are
\begin{equation}\tag{BCE21}
\begin{gathered}
\kappa(\epsilon)=3v_1,\quad
\kappa(T)=v_3-2v_2,\quad \kappa(z)=3v_2,\qquad F_H=0,\\
\Delta=\kappa F_M-F_H\kappa,\qquad
\Delta(T)=27cv_2,\quad\Delta(\epsilon)=\Delta(z)=0 .
\end{gathered}
\end{equation}
The free coefficients prove injectivity and the displayed defect
follows by applying \(F_M\) to \(T\).

There is a chain map
\(\mathcal K_m(M)\to\mathcal K_m(H)\) with its degree-zero
component exactly the original \(\kappa\) if and only if
\begin{equation}\tag{BCE22}
27c\in3^mO .
\end{equation}
Its degree-one component, when it exists, is unique and is
\begin{equation}\tag{BCE23}
\kappa_m^1=\kappa+3^{-m}\kappa F_M,\qquad
\kappa_m^1(T)=v_3-2v_2+3^{3-m}cv_2,\qquad
\kappa_m^1(\epsilon)=3v_1,\quad\kappa_m^1(z)=3v_2 .
\end{equation}
The chain equation is
\(\kappa_m^1(3^m-F_M)=3^m\kappa\).
Over \(K\), BCE10 forces BCE23, including its positive correction.
It is integral exactly under BCE22, and multiplication proves the
equation directly. Because the source and target are free, the
rational uniqueness also proves integral uniqueness.
For \(c\ne0\), BCE22 is \(m\le3+v_3(c)\); for \(c=0\) it holds
for every \(m\).

At larger \(m\), the exact obstruction to this fixed-degree-zero
lift is
\begin{equation}\tag{BCE24}
\operatorname{ob}_m(\kappa):M\longrightarrow H/3^mH,\qquad
T\longmapsto27cv_2\pmod{3^mH},\quad\epsilon,z\longmapsto0 .
\end{equation}
For \(c\ne0\) and \(m>3+v_3(c)\), its image has exact order
\(3^{m-3-v_3(c)}\). This follows by the valuation of its one
nonzero coefficient in the free \(v_2\) coordinate. It is the
obstruction to the specified lift, not to every possible
comparison between these objects.

Where BCE23 exists and \(m\ge2\), its induced degree-one map has
the evaluated composition
\begin{equation}\tag{BCE25}
(\kappa_m^1)_*\delta_m(1)
=[3^{m-1}\kappa_m^1(z)]
=[3^m v_2]=0\quad\hbox{in }H/3^mH .
\end{equation}
The nonzero class detected in BCE20 is therefore lost by this
particular cotangent receiver. The nonzero residue character
and the zero receiver are both retained.
The chain lift has two different components; it does not
remove the original coefficient Frobenius defect. Indeed
\(\kappa_m^1F_M(T)=27cv_2\) still, while \(F_H=0\).

\subsection{All-prime boundary families and the full support space}

Let \(B=B_{r,m,j}(O)\) with \(t=3^m\). By BCE3 it is free.
The analogue of BCE9 with \(O[t]\otimes_OB\) gives
\begin{equation}\tag{BCE26}
\begin{gathered}
\operatorname{RHom}_{O[t]}(B,E_F)
=\left[\operatorname{Hom}_O(B,E)
       \xrightarrow{\,3^m-F_E\,}\operatorname{Hom}_O(B,E)\right],\\
\operatorname{Ext}^1_{O[t]}(B,M_k)
 =\operatorname{Hom}_O(B,\operatorname{coker}L_{m,k}),\qquad
\operatorname{Hom}_{O[t]}(B,M_k)=0 .
\end{gathered}
\end{equation}
Injectivity of the resolution is checked by the highest polynomial
coefficient. Projectivity of \(B\) implies that
\(\operatorname{Hom}_O(B,-)\) is exact, proving the cokernel
identity. For an infinite basis this Hom is a product, not a
tensor direct sum; the formula retains that distinction.
Applying this exact functor to BCE13--25 gives their entire
all-prime family, not a finite set of tested characters.

In particular choose distinct \(A,D\in\mathcal A_{r,m}\), a prime
\(p\), and the balanced point with infinity coordinate \(A\),
\(p\)-coordinate \(D\), and every other coordinate \(A\).
For \(j=1\), define \(\ell:B_{r,m,1}(O)\to O\) by taking the
\(\vartheta_p\)-coefficient and evaluating it at this point.
The cylinder \(f=\mathbf1_{\{x_\infty=A,\ x_p=D\}}\) vanishes on
every endpoint and satisfies \(\ell(f\otimes\vartheta_p)=1\).
For \(m\ge2\), the input \(\ell\bmod9\in\operatorname{Hom}_O(B,Q)\)
therefore has connecting class
\begin{equation}\tag{BCE27}
b\longmapsto[3^{m-1}z\,\ell(b)],\qquad
\lambda_m\delta_m(\ell)(f\otimes\vartheta_p)=1/3\pmod O .
\end{equation}
For \(m=1\) every connecting class is zero, by BCE16.
The first nonzero count is \(m=2\), first present at \(r=3\);
one may take \(A=\{1,2\}\), \(D=\{1,3\}\).
Each rational prime supplies such a witness. This describes
the actual integral character module embedded by BCE4; it
does not claim a new Frobenius on the original Schwartz source.

The relation to the full support spectrum is also explicit.
For any finite nontrivial bounded distributive \(L\), let
\(Y=\operatorname{Spec}G_L(\mathbb Z)=S\triangleright
\operatorname{Spec}\mathbb Z\), as in FL1--15 and CFF28--35.
The exact constant-sheaf coupling
\(\mathcal U(E)=(c_XE,c_SE,\Delta)\) places the original
coefficient module on every stalk. Its ordinary global
cohomology is \(E[0]\); its supported cochain complex is
\(P_S\otimes_OE\), where
\[
P_S=\operatorname{Cone}(O[0]\xrightarrow{\Delta}C^\bullet(S,O))[-1],
\quad P_S^0=O,\quad P_S^i=C^{i-1}(S,O)\ (i\ge1),
\]
with differential \(-\Delta\) followed by \(-d\).
These are the exact proved CFF31 maps; all terms of \(P_S\)
are finite free. Hence applying \(\mathcal U\), or
\(P_S\otimes_O-\), to BCE13 gives respectively the stalkwise
extension or the short exact sequence of supported total
complexes. The total differential is
\(d_P\otimes1+(-1)^i1\otimes d_{\mathcal K}\) on bidegree
\((i,*)\); thus the two cone signs are fixed.
No nonvanishing in supported cohomology is inferred merely
from BCE27. For the two-element lattice,
\(P_S=[O\xrightarrow{-1}O]\) is contractible, and its tensor
complex is contractible as well: \(h_P\otimes1\) contracts
the displayed total differential, since the mixed terms
cancel with their opposite signs. Ordinary cohomology and
all stalks still retain BCE13--27.

Finally the separate labelled carrier is
\begin{equation}\tag{BCE28}
G_L(E)=\{(v,1_L):v\in E\}\cup\{(0,\lambda):\lambda\in L\},
\quad G_L(f)(v,\lambda)=(fv,\lambda).
\end{equation}
Addition uses the original join of labels. Evaluating the two
types of element proves that every additive map above lifts,
preserves composition, and fixes all zero labels. A nonzero
class sent to amplitude zero has image \(e=(0,1_L)\), while
\(\tau=(0,0_L)\) remains distinct. The torsion detector lifts
to \(G_L(K/O)\) with the same labels. No ordering or positivity
on this torsion target has been invented.

\subsection{The actual prism ideal and its Breuil--Kisin twist}

The comparison BCE8 is not the definition of syntomic cohomology.
Bhatt--Lurie's construction uses the absolute Nygaard filtration,
the Breuil--Kisin twist, and the difference between its partially
defined Frobenius and inclusion \cite{BCEBL}. We now calculate the
corresponding \emph{local twisted fixed-point complex at the given
prism}. This terminology specifies the calculation's domain:
we do not identify sections at one prism with the derived global
sections of the absolute prismatic site.

The bounded prism proved in DP1--16 is
\begin{equation}\tag{BCE29}
(C,I),\qquad I=(d),\qquad d=3+\epsilon,\qquad
d(3-\epsilon)=9,\qquad \phi(d)=3 .
\end{equation}
In particular \(d\) is a nonzerodivisor. The quotient map
\(q:C\to O\), \(q(\epsilon)=q(T)=0\), is a map of prisms to
\((O,(3))\): it commutes with Frobenius, hence with the
delta-operation \((\phi(x)-x^3)/3\), and sends \(I\) to \((3)\).
The nilpotent ideal \(M\) has \(M^4=0\), and
\(\phi^2|_M=0\). Neither equality identifies \(d\) with \(3\).

Write \(L=C\{1\}\) for the Breuil--Kisin line. We use the
following exact source facts: it is invertible and commutes
with base change of prisms; its semilinear Frobenius linearizes
to an isomorphism \(\phi^*L\simeq d^{-1}L\); and the crystalline
line \(O\{1\}\) has its canonical generator \(e_O\) with
\(\phi(e_O)=e_O/3\). These are the statements labelled
\texttt{remark:BK-functoriality},
\texttt{remark:Frobenius-on-twist}, and
\texttt{BKcrystalline} in the original Bhatt--Lurie source.

Choose a lift \(e\in L\) of \(e_O\). It generates \(L\):
\(C\) is local with nilpotent ideal \(M\) and quotient \(O\),
and a lift of a generator modulo this ideal is a generator.
Thus, for an actual unit \(u\in C^\times\),
\(\phi(e)=u e/d\). Reduction to \(O\) gives \(q(u)=1\).
Retain this initial generator, its unit, and the following
explicit comparison:
\begin{equation}\tag{BCE30}
v=u\phi(u),\qquad e_*=v e,\qquad
\phi(e_*)=e_*/d,\qquad
e=v^{-1}e_*,
\quad v^{-1}=\sum_{i=0}^{3}(-1)^i(v-1)^i .
\end{equation}
Indeed \(\phi^2(u)=1\), so \(\phi(v)=\phi(u)\) and
\(\phi(v)u=v\). The inverse formula follows from \(M^4=0\).
This comparison is independent of the chosen lift \(e\).
If \(e'=w e\) with \(q(w)=1\), its unit is
\(u'=\phi(w)u/w\), and
\[
u'\phi(u')w=u\phi(u)\phi^2(w)=u\phi(u).
\]
It is also the unique generator reducing to \(e_O\) and
satisfying the displayed Frobenius equation: the ratio of
two such generators satisfies \(\phi(w)=w\), hence
\(w=\phi^2(w)=q(w)=1\).
Thus BCE30 is a proved coordinate map retaining the original
unit, not an omission of it or of the distinguished element.

For \(m\ge0\), use \(e_*^{\otimes m}\) to display the line
\(C\{m\}\). Its Frobenius is exactly \(x\mapsto\phi(x)/d^m\).
Define
\begin{equation}\tag{BCE31}
N_mC=\{x\in C:\phi(x)\in d^mC\},\qquad
J_m(C)=[N_mC\xrightarrow{\ \phi/d^m-\iota\ } C]
\end{equation}
as a complex of \(O\)-modules in degrees zero and one.
The Frobenius difference is \(O\)-linear, not asserted
to be \(C\)-linear. Division means the unique preimage
under multiplication by the nonzerodivisor \(d^m\).
The sign is that of Bhatt--Lurie's syntomic fiber,
\(\phi\{m\}-\iota\). This defines a local complex without
assuming a descent theorem for it.

\subsection{Complete local filtration and its quotient defect}

For \(m\ge1\), let
\[
\eta_m=\begin{cases}1,&3\nmid m,\\0,&3\mid m,\end{cases}
\qquad
h_m=\max\{0,m-2-v_3(c)\},
\]
where \(h_m=0\) if \(c=0\). In the original ordered
coordinates \(1,\epsilon,T,z\) the full result is
\begin{equation}\tag{BCE32}
N_mC=
3^{m+\eta_m}O\,1\oplus O\epsilon
\oplus3^{h_m}OT\oplus Oz,\qquad
N_mM=O\epsilon\oplus3^{h_m}OT\oplus Oz .
\end{equation}
Here \(N_mM=N_mC\cap M\), and its divided Frobenius has
image in \(M\). To prove every coefficient, retain
\begin{equation}\tag{BCE33}
d^m=3^m+m3^{m-1}\epsilon .
\end{equation}
This is the full binomial expansion because \(\epsilon^2=0\).
More explicitly, for
\(x=x_0+x_\epsilon\epsilon+x_TT+x_zz\) one has
\[
\phi(x)=x_0+9cx_Tz,\qquad
d^{-m}\phi(x)
=3^{-m}x_0-m3^{-m-1}x_0\epsilon+3^{2-m}cx_Tz
\quad\text{in }C[1/3].
\]
All these coordinates are integral precisely when
\[
x_0\in3^mO,\quad mx_0\in3^{m+1}O,\quad
9cx_T\in3^mO .
\]
The first two conditions are \(x_0\in3^{m+\eta_m}O\);
the third is \(x_T\in3^{h_m}O\). No conditions are imposed
on \(x_\epsilon,x_z\), proving BCE32.

Write a general element of the domain in its retained basis as
\(3^{m+\eta_m}a+b\epsilon+3^{h_m}tT+fz\).
The complete differential is
\begin{equation}\tag{BCE34}
\begin{split}
(\phi/d^m-\iota)(x)
={}&3^{\eta_m}(1-3^m)a\\
&+(-m3^{\eta_m-1}a-b)\epsilon
-3^{h_m}tT+(3^{2+h_m-m}ct-f)z .
\end{split}
\end{equation}
The apparent negative powers are integral by the definitions
of \(\eta_m,h_m\). In particular \(m/3\in O\) when
\(\eta_m=0\). The factor \(1-3^m\) is a unit.
The constant coordinate in BCE34 first forces \(a=0\)
for an element of its kernel; then the \(T\) coordinate forces
\(t=0\), and the other coordinates force \(b=f=0\).
The independent columns for \(\epsilon,z\) kill those
coordinates in the cokernel. The two remaining relations
are exactly \(3^{\eta_m}(1-3^m)\) in the original constant
coordinate and \(-3^{h_m}\) in the original \(T\) coordinate.
Consequently
\begin{equation}\tag{BCE35}
\begin{gathered}
H^0J_m(C)=H^0J_m(M)=0,\\
H^1J_m(C)=(O/3^{\eta_m})[1]\oplus(O/3^{h_m})[T],
\qquad H^1J_m(M)=(O/3^{h_m})[T].
\end{gathered}
\end{equation}
These identifications are the coordinate maps
\([x]\mapsto(x_0\bmod3^{\eta_m},x_T\bmod3^{h_m})\);
the original off-diagonal terms remain in BCE34.
For \(m=0\), \(N_0C=C\) and \(\phi-I\) is zero on \(O\)
and invertible on \(M\), with inverse \(-I-\phi\).
Thus \(H^0J_0(C)=H^1J_0(C)=O\), while \(J_0(M)\)
is contractible.

The constant-coordinate torsion in BCE35 has an exact
geometric source: the filtration does not surject under
the original prism quotient \(q\). Indeed
\begin{equation}\tag{BCE36}
N_mO=3^mO,\qquad q(N_mC)=3^{m+\eta_m}O,\qquad
N_mO/q(N_mC)=3^mO/3^{m+\eta_m}O .
\end{equation}
The quotient complex \(J_m(C)/J_m(M)\) is
\[
\overline J_m=
[3^{m+\eta_m}O\xrightarrow{\,3^{-m}-1\,}O].
\]
There is an exact sequence of complexes
\begin{equation}\tag{BCE37}
0\longrightarrow\overline J_m\longrightarrow
J_m(O)\longrightarrow
[3^mO/3^{m+\eta_m}O\longrightarrow0]
\longrightarrow0 .
\end{equation}
Here \(J_m(O)\) is contractible since its differential,
on the generator \(3^m\), has value \(1-3^m\), a unit.
The connecting map takes the class of \(3^m\) to
\([1-3^m]\in H^1\overline J_m\), an isomorphism.
Thus the order-three constant class, present exactly when
\(3\nmid m\), is the calculated defect of filtered
surjectivity. It would be lost by replacing \(3+\epsilon\)
with \(3\) before doing the calculation.

\subsection{The exact return from the earlier collision comparison}

There is a chain map retaining both complexes and signs:
\begin{equation}\tag{BCE38}
\begin{array}{ccc}
M&\xrightarrow{\,3^m-\phi\,}&M\\
{\scriptstyle 3^m}\downarrow&&\downarrow{\scriptstyle -1}\\
N_mM&\xrightarrow{\,\phi/d^m-\iota\,}&M .
\end{array}
\end{equation}
For its domain, \(\phi(M)\subset Oz\) and
\(d^mz=3^mz\). Hence \(3^mM\subset N_mM\), and the
lower differential on \(3^mx\) is \(\phi(x)-3^mx\),
which proves the square. For \(m\ge1\), the map on
degree-one cohomology is the explicit surjection
\begin{equation}\tag{BCE39}
M/(3^m-\phi)M\longrightarrow O/3^{h_m},\qquad
[x]\longmapsto-x_T\bmod3^{h_m}.
\end{equation}
Well-definedness also follows directly from \(h_m\le m\)
and the fact that \(\phi(x)\) has no \(T\) coordinate.
In particular the previous nonzero order-three connecting
class has zero image, with an actual primitive:
\begin{equation}\tag{BCE40}
-3^{m-1}z
=(\phi/d^m-\iota)(3^{m-1}z),\qquad m\ge2 .
\end{equation}
The surviving \(T\) class when \(m>2+v_3(c)\), and the
constant class when \(3\nmid m\), have just been computed;
neither is substituted for the vanished \(z\) class.
Applying \(\operatorname{Hom}_O(B_{r,m,j}(O),-)\) preserves
this exact calculation since BCE3 proves that the source
is free. The chain map therefore kills the entire family
BCE27, not merely one chosen prime witness. Applying the
constant-sheaf coupling or the supported total-complex
construction following BCE27 preserves the same chain
map and primitive. The separate label functor BCE28
keeps every zero-support label on that map.

Finally we specify why BCE35 is not a computation of
absolute syntomic cohomology. The original prism quotient is
\begin{equation}\tag{BCE41}
D=C/(3+\epsilon)\cong(O/9)[T]/(T^2),\qquad
D/3D=\mathbb F_3[T]/(T^2).
\end{equation}
Frobenius on the second ring has image only the constants,
so is not surjective. A quotient of a perfectoid ring has
surjective Frobenius modulo \(3\), since roots lift from
the perfectoid ring. Thus \(D\) is not semiperfectoid,
and the quasiregular semiperfectoid formula for absolute
prismatic cohomology cannot be applied to identify it with
the given ring \(C\).
The actual absolute construction in \cite{BCEBL}, labelled
\texttt{construction:absolute-Nygaard-untwisted}, retains
derived global sections and a de Rham pullback; its syntomic
fiber is labelled \texttt{construction:syntomic-complex}.
Those terms have not been evaluated here. BCE30--40 instead
compute the local twist, domain, differential, and receiving
map completely. They show exactly which previous class is
lost before any proposed global arithmetic comparison.

\subsection{Completion retains the defect through its transported filtration}

The previous computation permits an exact completion comparison,
without identifying the result with absolute Nygaard completion.
Put \(R=\ker(\phi:C\to C)\). Directly from BCE5,
\begin{equation}\tag{BCE42}
\begin{gathered}
R=O\epsilon\oplus Oz\quad(c\ne0),\qquad R=M\quad(c=0),\\
\widehat C_N:=\varprojlim_m C/N_mC=C/R
=\begin{cases}O[T]/(T^2),&c\ne0,\\O,&c=0.\end{cases}
\end{gathered}
\end{equation}
Here equality with the quotient is the isomorphism induced by
the original quotient map \(q_N:C\to C/R\).
Indeed BCE32 gives \(\bigcap_mN_mC=R\). On \(C/R\),
the remaining coordinate exponents \(m+\eta_m\) and,
when present, \(h_m\), tend to infinity and differ from
\(m\) by bounded constants. Their filtration is cofinal
with the \(3\)-adic filtration of the finite free \(O\)-module
\(C/R\), which is complete. This proves both kernel and
surjectivity of the inverse-limit map. The quotient tower
has surjective transition maps, so its first derived
inverse limit is zero as well. The original multiplication
descends because \(R\) is an ideal; \(T^2=6\epsilon\)
becomes \(T^2=0\), with its quotient map recorded.

Write \(\bar C=C/R\). The transported filtration and
induced Frobenius are exactly
\begin{equation}\tag{BCE43}
\begin{gathered}
\bar N_m=q_N(N_mC)
=3^{m+\eta_m}O\,1\oplus3^{h_m}OT,\quad
\bar\phi(x_0+x_TT)=x_0,\quad q_N(d)=3,\\
\widetilde N_m:=\{x\in\bar C:\bar\phi(x)\in3^m\bar C\}
=3^mO\,1\oplus OT .
\end{gathered}
\end{equation}
Every \(T\)-summand in BCE43 is omitted when \(c=0\).
Thus recomputing a Frobenius preimage after taking this
quotient is a different operation from transporting the
original filtration. Their domains agree exactly when
\(\eta_m=h_m=0\). For \(c\ne0\), this means
\(3\mid m\) and \(m\le2+v_3(c)\); for \(c=0\), it
means \(3\mid m\). In every other case the inclusion
\(\bar N_m\subset\widetilde N_m\) is strict.

Define
\(\bar J_m=[\bar N_m\xrightarrow{\bar\phi/3^m-\iota}\bar C]\)
and
\(\widetilde J_m=[\widetilde N_m
\xrightarrow{\bar\phi/3^m-\iota}\bar C]\).
These remain complexes of \(O\)-modules; no quotient
\(\bar C\)-module structure on their cohomology is presumed.
The divided maps descend from \(J_m(C)\): \(\phi(R)=0\),
and applying \(q_N\) to \(d^my=\phi(x)\) gives
\(3^mq_N(y)=\bar\phi(q_N(x))\).
Both quotient source and quotient target are torsionfree,
so this specifies the divided map uniquely.
The kernel of \(J_m(C)\to\bar J_m\) is exactly
\([R\xrightarrow{-1}R]\), proving that the quotient is
a quasi-isomorphism.

For \(m\ge1\), the other complex \(\widetilde J_m\)
is contractible: its differential sends \(3^m\) to the
unit \(1-3^m\), and \(T\), when present, to \(-T\).
The precise difference between the two complexes is
\begin{equation}\tag{BCE44}
\begin{gathered}
Q_m=\widetilde N_m/\bar N_m
=(O/3^{\eta_m})[3^m]\oplus(O/3^{h_m})[T],\\
0\longrightarrow\bar J_m\longrightarrow\widetilde J_m
\longrightarrow[Q_m\longrightarrow0]\longrightarrow0,\\
Q_m\xrightarrow[\sim]{\ \partial\ }H^1\bar J_m,\qquad
[3^m]\longmapsto[1-3^m],\quad[T]\longmapsto[-T].
\end{gathered}
\end{equation}
Again omit \(T\) when \(c=0\).
The exact sequence is termwise exact by its definitions.
Lifting each displayed degree-zero quotient generator and
applying its original differential gives the connecting
formula and its signs. Contractibility of
\(\widetilde J_m\) proves the isomorphism. This recovers
every group in BCE35, including both summands, as the
explicit defect of replacing the transported filtration.
Although \(\epsilon\) vanishes under completion, its
\(\eta_m\)-contribution remains in that filtration.
Thus the completed ring and endomorphism alone do not
retain all the local comparison data.

There is also an explicit return and chain homotopy for
the completion map \(\pi:J_m(C)\to\bar J_m\), for \(m\ge1\).
Let \(s^1:\bar C\to C\) lift the original coordinates \(1,T\)
identically, omitting \(T\) when \(c=0\). For a domain element
\(\bar x=3^{m+\eta_m}a+3^{h_m}tT\), put
\begin{equation}\tag{BCE45}
\begin{gathered}
s^0(\bar x)=3^{m+\eta_m}a+3^{h_m}tT
-m3^{\eta_m-1}a\epsilon+3^{2+h_m-m}ctz,\\
P_R(y)=y-s^1\pi(y),\qquad
h(y)=-P_R(y):J_m(C)^1\longrightarrow J_m(C)^0,\\
\pi s=I,\qquad d_Jh+hd_J=I-s\pi .
\end{gathered}
\end{equation}
When \(c=0\), omit both terms containing \(t\).
The added terms of \(s^0\) lie in \(R\subset N_mC\).
Substituting into BCE34 cancels its two off-diagonal
entries, proving \(d_Js^0=s^1d_{\bar J}\).
In degree one, \(d_Jh(y)=P_R(y)\), since
\(d_J|_R=-I\). In degree zero, writing \(x\) in the
basis of BCE34 gives
\[
-P_R(d_Jx)=x-s^0\pi(x),
\]
with the same two displayed corrections; for \(c=0\)
the \(T\)-coordinate lies in \(R\) and obeys the same
identity. These calculations prove the two homotopy
identities on every original coordinate.
The return is \(O\)-linear, not a section of rings.
It preserves the local cohomology without pretending
that the completed coefficient ring retains the
original multiplication or the untransported filtration.

\subsection{Exact conclusions and the remaining arithmetic comparison}

The computation closes the coefficient-extension calculation
for the actual balanced characters: the first mixed count
kills the connecting map, the second retains an order-three
class, and BCE19 detects it exactly. The old cotangent
injection either has the unique integral chain repair BCE23,
which kills this class, or has the evaluated obstruction BCE24.
All tensor comparison groups and their perfect torsion
pairings are given in BCE11--12.

These results do not impose positivity on the original zeta
pairing. Their characteristic-zero comparison complexes are
contractible by BCE10, while the original complex boundary
module of BCE2 itself remains nonzero. Those are different
objects joined by the explicit coefficient maps BCE3--4,
not a license to discard the boundary. On the universal
\(\mathbb Z_{(3)}[c]\) presentation, extension to
\(\mathbb C[c]\) gives the same inverse BCE10 since 3 is
invertible; extension to \(O\) retains the computed torsion.
The classical zeta function, its poles, Gamma factors,
completion multiplier and original Weil form have not been
replaced or altered by any of these coefficient operations.
An arithmetic sign comparison must act on that form and
account for these exact zero and nonzero images. No such
sign conclusion follows from this calculation alone.
Moreover, using the actual prism ideal and its twist changes
the receiving complex as in BCE31--40: the earlier order-three
class has an explicit primitive there, while different
filtration defects remain. An absolute or arithmetic
interpretation must retain that calculated distinction.

\input{collision_boundary/BCE_FIGURE.tex}
